\section{A necessity theorem for causal realization}\label{sec:v8-congruence}
We now characterize the causal constraint without imposing a tree
geometry. The time dependence of a machine is retained, so the result
does not replace a free clock by an uncharged stationary-state assumption.
Fix a horizon $T$, finite report alphabets, one Bayesian law, and
square-integrable Hilbert targets $U_t$. There is no control in this
section. Let $\mathcal H_t$ be the finite set of positive-probability
histories, $p_t(h)=\Pp(H_t=h)$, and
$m_t(h)=\E[U_t\mid H_t=h]$. Let $\beta_t\ge0$ be fixed weights.

A family of partitions $(\mathcal P_t)_{t=0}^T$, with
$\mathcal P_0=\{\emptyw\}$, is \emph{successor compatible} if
\begin{equation}
 h,h'\in C\in\mathcal P_t,\quad hy,h'y\in\mathcal H_{t+1}
 \quad\Longrightarrow\quad
 hy,h'y\text{ belong to the same cell of }\mathcal P_{t+1}.
                                                               \label{eq:v8-congruence}
\end{equation}
Only common positive-probability continuations enter this condition.
For $C\in\mathcal P_t$ put
$\overline m_C=p_t(C)^{-1}\sum_{h\in C}p_t(h)m_t(h)$ and define
\begin{equation}
 \mathcal V(\mathcal P)=
 \sum_{t=1}^T\beta_t\sum_{C\in\mathcal P_t}
              \sum_{h\in C}p_t(h)\|m_t(h)-\overline m_C\|^2.
                                                               \label{eq:v8-variancefunctional}
\end{equation}

\begin{theorem}[Recursive partition characterization]
\label{thm:v8-congruence}
For the cardinality signature, allowing randomized time-dependent
updates and decoders and independent public randomness, the optimal
weighted excess over the full-history Bayes risk equals
\begin{equation}
 \inf_{\#I_t\le S}
 \sum_{t=1}^T\beta_t\E\|m_t(H_t)-\widehat U_t\|^2
 =\min_{\substack{\mathcal P\ \mathrm{successor\ compatible}\\
                         |\mathcal P_t|\le S}}
                         \mathcal V(\mathcal P).       \label{eq:v8-exact-congruence}
\end{equation}
Thus same-order realization of a sequence of static checkpoint optima
is equivalent to the existence of successor-compatible partitions
with the corresponding widths and variance bounds. This is a necessary
as well as a sufficient condition. It is invariant under relabelling of
histories; suffix closure is one sufficient construction of these
partitions, not a necessary choice of coordinates.
\end{theorem}
\begin{proof}
A deterministic initialized machine partitions $\mathcal H_t$ according
to $I_t$. Its update depends only on $(t,I_t,y)$, so its partitions satisfy
\eqref{eq:v8-congruence}. Given a state cell, the mean $\overline m_C$
minimizes its squared error by Hilbert orthogonality. The resulting
excess is precisely \eqref{eq:v8-variancefunctional}.

Conversely, label each partition's cells by at most $S$ labels. For each
cell $C$ and symbol $y$, send its label to the cell containing $hy$ for
any supported $h\in C$ with supported extension $hy$. Compatibility
makes the answer independent of the choice of $h$. If there is no such
extension, define the update arbitrarily. At time $t$ decode the label
as $\overline m_C$. Induction identifies the actual label with the cell
of the actual history. There are finitely many partition families, so
the minimum exists.

To include randomization, realize the finitely many update and decoder
kernels by independent uniforms, and condition on their entire tape.
For each tape, the rule is a deterministic time-dependent $S$-state
machine. Independence of the tape and the experiment preserves the
history law and conditional means. The average of its weighted costs
cannot be smaller than the deterministic minimum. Public independent
randomness is handled in the same way. This freezing argument concerns
one fixed Bayesian sum; it does not interchange a minimax task supremum
with the average over tapes.

Unrestricted checkpoint encoders correspond to arbitrary partitions at
each time, without \eqref{eq:v8-congruence}. Hence they have cost
$\sum_t\beta_t e_S^2(m_t(H_t))$. Formula
\eqref{eq:v8-exact-congruence} gives exactly the asserted equivalence of
order comparisons, including the case when a static cost is zero.
\end{proof}

\begin{corollary}[Finite-command cell instruments]
\label{cor:v8-a1}
Let $k_1,\ldots,k_J$ be Borel functions on a parameter space with
$k_j\ge0$ and $\sum_jk_j=1$, and fix a prior $\mu$. At each time the
raw cell is drawn independently conditional on the parameter. A command
$g$ from a prescribed finite menu in $[\eta,1-\eta]^J$, $0<\eta<1/2$,
has a prescribed parameter-independent law $\nu_t$. Cell $j$ is
reported as $j$ with probability $g_j$ and otherwise as the single
failure symbol $\dagger$. Only the command and the reported symbol are
observed. Fix a finite horizon, finite future-query menus with known
likelihoods $H_{tq}\in[0,1]$, independent query weights $\omega_{tq}$,
and nonnegative checkpoint weights $\beta_t$.

For the persistent-cardinality signature, the optimal common causal
encoder for the weighted sum of quadratic prediction excesses is exactly
\eqref{eq:v8-exact-congruence}, with the finite history law and query
means computed from the actual cell likelihoods. This applies to the
finite-command specialization of the A1 finite-cell protocol
\cite{QiA1v37}. It is an exact scalarized criterion, not an asserted
solution of the different maximum-over-checkpoints criterion with
private randomization.
\end{corollary}
\begin{proof}
For the observed pair $(g,x)$ define
\[
 \lambda_{g,x}(\theta)=
 \begin{cases}
 g_x k_x(\theta),&x\in\{1,\ldots,J\},\\
 \sum_j(1-g_j)k_j(\theta),&x=\dagger.
 \end{cases}
\]
For a history $h=((g_s,x_s))_{s\le t}$ its probability is
$\prod_{s\le t}\nu_s(g_s)$ times
$\int\prod_{s\le t}\lambda_{g_s,x_s}\,d\mu$. On positive-probability
histories the mean of query $q$ is
\[
 m_{tq}(h)=
 \frac{\int H_{tq}(\theta)\prod_{s\le t}\lambda_{g_s,x_s}(\theta)
                    \,\mu(d\theta)}
      {\int\prod_{s\le t}\lambda_{g_s,x_s}(\theta)\,\mu(d\theta)}.
\]
Take the Hilbert prediction vector with coordinates
$\sqrt{\omega_{tq}}\,m_{tq}(h)$ and apply
Theorem~\ref{thm:v8-congruence}. All histories are finite, and the law
used there is exactly the actual acquisition law above. A failure
history contains $\dagger$, not the unobserved cell $j$; the successor
compatibility condition therefore requires one common update for that
report. Integration of the given Borel likelihoods may require arbitrary
real constants or offline work. Neither is claimed finite under the
cardinality signature. Zero-evidence histories are omitted rather than
given artificial posteriors.
\end{proof}

\begin{corollary}[Continuation obstruction]\label{cor:v8-continuation}
Let $g$ be a bounded readout on all finite words over a finite alphabet.
For stationary deterministic machines required to reproduce $g$ on every
word, an exact $S$-state realization exists if and only if the relation
\begin{equation}
 h\equiv_g h'\quad\Longleftrightarrow\quad
                g(hw)=g(h'w)\text{ for every finite }w
                                                               \label{eq:v8-nerode}
\end{equation}
has at most $S$ classes. More generally put
$d_*(h,h')=\sup_w\|g(hw)-g(h'w)\|$. Every stationary deterministic
realization with uniform error at most $\epsilon$ has state fibres of
$d_*$-diameter at most $2\epsilon$, and thus at least as many states as
any $2\epsilon$-separated set of histories.
\end{corollary}
\begin{proof}
Histories in one machine state have the same state after every common
continuation. Exact output equality implies \eqref{eq:v8-nerode}; for
approximate outputs the triangle inequality gives the diameter bound.
The relation \eqref{eq:v8-nerode} is a right congruence. Its classes,
updated by appending the new symbol, form an exact machine, whose output
is $g(h)$, well defined by the empty continuation. This is the output
version of the classical finite-automaton congruence argument
\cite{Nerode1958}; the finite-horizon, weighted-distortion characterization
in Theorem~\ref{thm:v8-congruence} is the form needed here.
\end{proof}

\begin{example}[A two-valued checkpoint requiring exponentially many
persistent states]\label{ex:v8-index}
At the first time observe a uniformly distributed word
$X\in\{0,1\}^d$. At the second time observe an independent uniform index
$J\in\{1,\ldots,d\}$ and predict $U=X_J$. The raw first acquisition costs
$d$ bits; it is not called a one-bit observation. The full-history
conditional mean at the second checkpoint has only two values, so its
static two-label error is zero. Nevertheless a causal exact predictor
requires $S\ge2^d$. Two different first-stage words in one state differ
in a coordinate which can be queried with positive probability. The
second-stage update cannot separate them. Retaining the whole word is
sufficient.

There is also a quantitative obstruction. If $e$ is the smallest error
probability of an $S$-state predictor, $2\le S\le2^d$, then
\begin{equation}
 e\ge h_2^{-1}\!\left(1-\frac{\log_2 S}{d}\right),     \label{eq:v8-index-fano}
\end{equation}
where $h_2$ is the binary entropy and its inverse takes values in
$[0,1/2]$. Indeed, condition also on every independent public coding
variable $R$. Since the first register has at most $S$ values,
$H(X\mid I_1,R)\ge d-\log_2 S$. Predicting each coordinate optimally
from $(I_1,R)$ gives errors $e_j\le1/2$, and the binary entropy inequality
and concavity give
\[
 H(X\mid I_1,R)\le\sum_{j=1}^d H(X_j\mid I_1,R)
             \le\sum_{j=1}^d h_2(e_j)
             \le d\,h_2\!\left(d^{-1}\sum_j e_j\right).
\]
The average on the right is a lower bound for any actual queried-bit
error. This proves \eqref{eq:v8-index-fano}. If squared loss is used
instead, its optimal value is at least $e/2$, since
$p(1-p)\ge\min(p,1-p)/2$. Thus the discrepancy persists for the
Hilbert-loss slice of the spectrum. It is a temporal obstruction, not
poor scalar quantization. The query must be answered after the update;
$J$ is not an additional persistent decoder input.
\end{example}
